\documentclass[../../../include/open-logic-section]{subfiles}

\begin{document}

\olfileid{sth}{cardinals}{hp}
\olsection{پیوست: اصل هیوم}

در~\olref[cp]{sec} اصل کانتور را توصیف کردیم. آن اصل چنین بود:
\begin{align*}
	\card{A} = \card{B} & \text{ اگر و تنها اگر } A \approx B.
\intertext{این اصل بسیار شبیه چیزی است که امروزه \emph{اصل هیوم}
نامیده می‌شود و می‌گوید:}
	\fregenum{x} {F(x)} = \fregenum{x}{G(x)} & \text{ اگر و تنها اگر } F \sim G
\end{align*}
که در آن «$F\sim G$» مخفف این است که تعداد~$F$ها دقیقاً با تعداد~$G$ها
برابر است؛ یعنی می‌توان~$F$ها را در تناظر دوسویه با~$G$ها قرار داد؛
یعنی:
\begin{align*}
	\exists R(&\forall v\forall y(Rvy \lif (Fv \land Gy)) \land {}\\
		&\forall v(Fv \lif \lexists![y][Rvy]) \land {}\\
		&\forall y(Gy \lif \lexists![v][Rvy]))
\end{align*}
اما میان اصل هیوم و اصل کانتور تفاوتی نوعی وجود دارد. در بیان اصل کانتور،
متغیرهای «$A$» و «$B$» ترم‌های مرتبهٔ اولی‌اند که نمایندهٔ
\emph{مجموعه‌ها} هستند. در بیان اصل هیوم، «$F$»، «$G$» و «$R$»
ترم‌های مرتبهٔ اول \emph{نیستند}؛ بلکه در \emph{جایگاه محمول} قرار
دارند. (شاید نمایندهٔ \emph{ویژگی‌ها} باشند.) بنابراین می‌توانیم اصل
هیوم را چنین شرح دهیم: عددِ~$F$ها با عددِ~$G$ها برابر است اگر و تنها اگر
$F$ها با~$G$ها تناظر دوسویه داشته باشند. این اصل «اصل هیوم» نامیده
می‌شود، زیرا هیوم زمانی چنین نوشت:
\begin{quote}
  هرگاه دو عدد چنان با هم ترکیب شوند که هر واحدِ یکی همواره واحدی متناظر
  در دیگری داشته باشد، آن‌ها را برابر می‌خوانیم.
  \citep[Pt.III Bk.1 \S1]{Hume1740}
\end{quote}
\citet[\S63]{Frege1884} این قطعه از هیوم را نقل کرد و سپس، در واقع،
طرحی از چیزی را که اصل هیوم نامیده‌ایم به دست داد؛ و بدین‌ترتیب اصل هیوم
را در کانون توجه معاصرِ منطق ریاضی قرار داد.

باید به شباهت ساختاری میان اصل هیوم و قانون پایهٔ پنجم توجه کنید. این
قانون را در~\olref[story][blv]{sec} چنین صورت‌بندی کردیم:
\[
	\fregeext{x}{F(x)} = \fregeext{x}{G(x)} \text{ اگر و تنها اگر } \lforall[x][(F(x) \liff G(x))].
\]
در اینجا قدری توضیح و مقایسه سودمند است.

اصلی مانند اصل هیوم یا قانون پایهٔ پنجم را به دو شیوه می‌توان فهمید:
\emph{محمولی} یا \emph{غیرمحمولی}
(\olref[story][predicative]{sec} را به یاد آورید). در قرائت غیرمحمولیِ
قانون پایهٔ پنجم، برای هر~$F$، شیء~$\fregeext{x}{F(x)}$ در همان دامنهٔ
سوربندی قرار می‌گیرد که برای صورت‌بندی خودِ قانون پایهٔ پنجم به کار
بردیم. به همین ترتیب، در قرائت غیرمحمولیِ اصل هیوم، برای هر~$F$، شیء
$\fregenum{x}{F(x)}$ در همان دامنهٔ سوربندی قرار می‌گیرد که برای
صورت‌بندی اصل هیوم به کار بردیم. در مقابل، بنا بر فهم \emph{محمولی}،
اشیای~$\fregeext{x}{F(x)}$ و~$\fregenum{x}{F(x)}$ موجوداتی از دامنه‌ای
\emph{دیگر} خواهند بود.

حال اگر قانون پایهٔ پنجم را غیرمحمولی بخوانیم، از راه فهم خام به
ناسازگاری می‌انجامد (برای جزئیات بنگرید به~\olref[story][blv]{sec}).
این قانون نیز، بسیار شبیه فهم خام، با قرائت \emph{محمولی} سازگارشدنی
است؛ اما احتمالاً همهٔ کارهایی را که از آن می‌خواستیم انجام نخواهد داد.

بااین‌حال، اصل هیوم را \emph{می‌توان} به‌طور سازگار غیرمحمولی خواند؛ و
در این قرائت بسیار نیرومند است.

برای نمونه، محمول «$x\neq x$» را در نظر بگیرید که آشکارا هیچ چیز آن را
ارضا نمی‌کند. اصل هیوم اکنون شیء~$\# x(x\neq x)$ را به دست می‌دهد.
می‌توانیم این شیء را عدد~$0$ بدانیم. اکنون بنا بر فهم
\emph{غیرمحمولی}---و \emph{فقط} بنا بر فهم غیرمحمولی---این موجود~$0$ در
دامنهٔ آغازین سوربندی ما قرار می‌گیرد. پس می‌توانیم اصل هیوم را به‌طور
معنادار با محمول «$x=0$» به کار بریم و شیء~$\#x(x=0)$ را به دست آوریم.
می‌توانیم این شیء را عدد~$1$ بدانیم. افزون بر این، اصل هیوم مستلزم
$0\neq1$ است، زیرا هیچ تناظر دوسویه‌ای میان اشیای نابرابر با خود و اشیای
برابر با~$0$ وجود ندارد (از گروه نخست هیچ شیئی و از گروه دوم یک شیء وجود
دارد). بار دیگر با کار غیرمحمولی، $1$ در دامنهٔ آغازین سوربندی ما قرار
می‌گیرد. پس می‌توانیم اصل هیوم را به‌طور معنادار با محمول
«$(x=0\lor x=1)$» به کار بریم و شیء~$\#x(x=0\lor x=1)$ را به دست
آوریم. می‌توانیم آن را عدد~$2$ بدانیم و نشان دهیم~$0\neq2$ و~$1\neq2$
و به همین ترتیب ادامه دهیم.

خلاصه آنکه اصل هیوم با قرائت غیرمحمولی مستلزم وجود \emph{بی‌نهایت شیء}
است. همین امر \emph{منطق‌گرایان نوفرگه‌ای} را تشویق کرده است اصل هیوم
را بنیاد حساب بگیرند.

اما خودِ فرگه اصل هیوم را بنیاد حساب خود قرار نداد. او اصل هیوم را از
تعریفی صریح اثبات کرد: $\fregenum{x}{F(x)}$ به‌منزلهٔ مصداق مفهوم
$F\sim\Phi$ تعریف می‌شود. با اصطلاحات امروزی شاید بکوشیم آن را چنین
بنویسیم: $\fregenum{x}{F(x)}=\Setabs{G}{F\sim G}$؛ اما این کار ما را
دوباره به مشکلات فهم خام بازمی‌گرداند.

\end{document}
